\documentclass[11pt]{article}
\usepackage{amsmath, amssymb}
\usepackage{geometry}
\geometry{margin=1in}
\usepackage{fancyhdr}
\usepackage{lastpage}
\usepackage{parskip} % Better paragraph spacing
\usepackage{titlesec} % Better section formatting
\usepackage{graphicx}
\usepackage{float}

% Define commands for reusable content
\newcommand{\coursename}{CS-401 Applied Data Analysis}
\newcommand{\lecturenum}{01}
\newcommand{\lecturedate}{\today}

% Configure section spacing
\titlespacing*{\section}{0pt}{2ex plus 1ex minus .2ex}{1.5ex plus .2ex}
% Header / Footer
\pagestyle{fancy}
\fancyhf{}
\fancyhead[L]{\coursename}
\fancyhead[R]{Lecture \lecturenum}
\fancyfoot[L]{EPFL - Computer Science}
\fancyfoot[C]{Author: Jonathan Dutoit}
\fancyfoot[R]{Page \thepage\ of \pageref{LastPage}}

\begin{document}

\begin{center}
  {\large EPFL - Computer Science} \\[0.5em]
  {\Large \textbf{\coursename}} \\[0.5em]
  {\large Lecture \lecturenum\ --- \lecturedate} \\[0.5em]
  {\large Author: Jonathan Dutoit} \\
\end{center}

\vspace{2cm}

\begin{center}
  {\LARGE \textbf{Introduction to ADA}}
\end{center}

\vspace{1cm}

\section*{Data analysis}
“... the process of inspecting, cleaning, transforming, and modeling data with the goal of discovering useful information, suggesting conclusions, and supporting decision-making.”

“Data analysis has multiple facets and approaches, encompassing diverse techniques under a variety of names, in different business, science, and social science domains.”

\section*{Applied Data Analysis}
\begin{itemize}
	\item This course is about \textbf{breadth}, not depth
	\item "What methods, principles, and tools are out there?, rather than "How can I become an expert in deep learning for computer vision applied to images of cats?"
	\item Data science is a fast-paced, shifting fields
	\item Obsessing on one tool or technique won't pay off in a few years
	\item Be ready to explore and keep learning on your own
\end{itemize}

\textbf{Goal:} Enable you to conduct a full-fledge data science project from start to finish

\section*{Syllabus}
\begin{itemize}
	\item Handling data
	\begin{itemize}
		\item "Slicing and dicing": obtaining, preparing, juggling data
	\end{itemize}	
	\item Visualizing data
	\begin{itemize}
		\item Exploration of data, communication of results
	\end{itemize}
	\item Describing data
	\begin{itemize}
		\item How to support (and be suspicious of) claims about data
	\end{itemize}
        \item Regression analysis for disentangling data
	\begin{itemize}
		\item How to disentangle datasets with correlated variables
	\end{itemize}
	\item Casual analysis of observational data
	\begin{itemize}
		\item How to deal with "found data"
		\item correlation != causation
	\end{itemize}
\end{itemize}

\section*{Syllabus (cont'd)}
\begin{itemize}
	\item Learning from data
	\begin{itemize}
		\item supervised learning
		\item unsupervised learning
		\item applied aspects of machine learning
	\end{itemize}
	\item Handling specific types of data
	\begin{itemize}
		\item handling text data
		\item handling network data
	\end{itemize}
	\item Scaling to massive data
\end{itemize}

\section*{Grading}

\begin{itemize}
	\item 15\% Homework assignment
	\item 50\% Final Exam
	\item 35\% Project 
\end{itemize}


\end{document}
